% \iffalse meta-comment
%
% hit-engine.dtx 是引擎层
%
% \fi
%
% \section{引擎层}
%
% \changes{v3.2a}{2026/09/20}{新增引擎层，模板 \hologo{XeLaTeX}、\hologo{LuaLaTeX} 与 \hologo{pdfLaTeX} 都能编，
%   报错文案里推荐引擎 \texttt{lualatex} 排在前面；\hologo{pdfLaTeX} 下中西文之间的间距在交界处自己补
%   （\pkg{CJK} 没有这个机制），\texttt{\textasciitilde} 还原成 \cs{nobreakspace}（\pkg{CJK} 的那个不健壮）；\cs{CJKglue} 的直接改写、
%   \cs{hit@engine@cjk@char:n}、\cs{hit@engine@latin@chars:n} 收进这里；
%   东八区的今天（\cs{c\_sys\_year\_int} 跟编译机器
%   时区走，国外编译封面日期差一天）；下划线按 Word 画：线心在基线下 0.135\,em、粗 0.05\,em}
%
% 中文排版的活儿是 \pkg{ctex} 干的，而 \pkg{ctex} 两个引擎都支持：\hologo{XeLaTeX} 底下它调 \pkg{xeCJK}，
% \hologo{LuaLaTeX} 底下调 \pkg{luatexja}。挡在门口的是模板自己直接用的那几个
% \pkg{xeCJK} 专属接口——它们在 \pkg{luatexja} 底下没有同名的东西。
%
% 本文件把这几处收成一层薄封装：每个功能给一个 \cs{hit@engine@} 开头的名字，
% 底下按引擎分两套实现，调用方一处 \cs{sys\_if\_engine} 都不用写。加引擎只动这里。
%
% \textbf{排在装配表第三位。} 要早于 \file{hit-fonts.dtx}（第五位，字体族里要用
% 字符归类）、\file{hit-footnote.dtx}（字形探测）、\file{hit-spread.dtx}（定宽拉伸）
% 与 \file{docxlike-format.dtx}（字距）。
%
%    \begin{macrocode}
%<*shared-engine>
%    \end{macrocode}
%
% \subsection{引擎守卫}
%
% 引擎是先决条件，不对就不必再往下走，所以守卫在装配表里排得比选项模块还早。
% \hologo{LuaLaTeX} 是主线，\hologo{XeLaTeX} 次之；\hologo{pdfLaTeX} 也能编：中文走
% \pkg{ctex} 的 \pkg{CJK} 后端，字体见 \file{src/engine/hit-fonts.dtx} 的“\hologo{pdfLaTeX} 的字体”，
% 只认直接出 PDF 的 \texttt{pdflatex}，\texttt{latex} 出 DVI 那条路挡在门外。
%    \begin{macrocode}
\bool_lazy_any:nF
  {
    { \sys_if_engine_xetex_p: }
    { \sys_if_engine_luatex_p: }
    { \sys_if_engine_pdftex_p: }
  }
  { \ClassError { hithesis } { Please~use:~\MessageBreak~lualatex~(or~xelatex,~pdflatex) } { } }
\sys_if_engine_pdftex:T
  {
    \int_compare:nNnT { \tex_pdfoutput:D } < { 1 }
      { \ClassError { hithesis } { Please~use~pdflatex,~not~latex~(DVI~output) } { } }
  }
%    \end{macrocode}
%
% \subsection{把码位归成汉字}
%
% \cs{hit@engine@cjk@chars:n}\marg{码位列表} 把列出的码位当汉字排，让它们从中文
% 字体取字形。用处见 \file{src/engine/hit-fonts.dtx}：\texttt{□} 这类几何符号按
% Unicode 区块归了西文，西文字体里没有就整个字排丢了。
%
% 两个引擎的分法不一样。\pkg{xeCJK} 是逐个码位声明字符类，一个码位一句；
% \pkg{luatexja} 是先把一批码位定义成一个「字符范围」，再把整个范围划归和文，
% 两句管一批。范围号 \pkg{luatexja} 预定义了 1--8、留给用户 1--217，这里从 100
% 起编，离预定义的那几个远一点；每调用一次占一个号。
%    \begin{macrocode}
\sys_if_engine_luatex:TF
  {
    \int_new:N \g__hit_engine_range_int
    \int_gset:Nn \g__hit_engine_range_int { 100 }
    \cs_new_protected:Npn \hit@engine@cjk@chars:n #1
      {
        \int_gincr:N \g__hit_engine_range_int
        \exp_args:Nx \ltjdefcharrange
          { \int_use:N \g__hit_engine_range_int } {#1}
        \exp_args:Nx \ltjsetparameter
          { jacharrange = { + \int_use:N \g__hit_engine_range_int } }
      }
  }
  {
    \cs_new_protected:Npn \hit@engine@cjk@chars:n #1
      { \clist_map_inline:nn {#1} { \xeCJKDeclareCharClass { CJK } {##1} } }
  }
%    \end{macrocode}
%
% \cs{hit@engine@latin@chars:n}\marg{码位列表} 反过来：把列出的码位交回西文，
% 让它们从西文字体取字形。破折号要用它，见 \file{src/engine/hit-fonts.dtx}。
%
% \pkg{xeCJK} 那边把字符类改成 \texttt{Default}——那一类的意思就是「\pkg{xeCJK}
% 不管」，字形跟当前字体走。\pkg{luatexja} 那边把码位划成一个范围再从和文里减掉，
% 范围号与上面同一条序列。
%    \begin{macrocode}
\sys_if_engine_luatex:TF
  {
    \cs_new_protected:Npn \hit@engine@latin@chars:n #1
      {
        \int_gincr:N \g__hit_engine_range_int
        \exp_args:Nx \ltjdefcharrange
          { \int_use:N \g__hit_engine_range_int } {#1}
        \exp_args:Nx \ltjsetparameter
          { jacharrange = { - \int_use:N \g__hit_engine_range_int } }
      }
  }
  {
    \cs_new_protected:Npn \hit@engine@latin@chars:n #1
      { \clist_map_inline:nn {#1} { \xeCJKDeclareCharClass { Default } {##1} } }
  }
%    \end{macrocode}
%
% \hologo{pdfTeX} 下 \pkg{CJKutf8} 把 0x80 以上的字节全接走了，码位不必再归成汉字；归成西文
% 只用得上破折号那一个，写成 \pkg{inputenc} 的 \cs{textemdash}。
%    \begin{macrocode}
\sys_if_engine_pdftex:T
  {
    \cs_set_protected:Npn \hit@engine@cjk@chars:n #1 { }
    \cs_set_protected:Npn \hit@engine@latin@chars:n #1
      {
        \clist_map_inline:nn {#1}
          {
            \int_compare:nNnT {##1} = { "2014 }
              { \DeclareUnicodeCharacter { 2014 } { \textemdash } }
          }
      }
  }
%    \end{macrocode}
%
% \subsection{汉字之间的字距}
%
% \cs{hit@engine@cjkglue@zero:} 把汉字之间的那段胶归零，
% \cs{hit@engine@cjkglue@set:nn}\marg{汉字间}\marg{中西文间} 给它们设值。
% 两个都是局部的，写在组里出组即还原。
%
% \pkg{xeCJK} 把这两段胶做成两个宏 \cs{CJKglue} 与 \cs{CJKecglue}，改它们等于改
% 排版时插进去的东西；\pkg{luatexja} 做成两个参数 \texttt{kanjiskip} 与
% \texttt{xkanjiskip}，走 \cs{ltjsetparameter}。名字不同，管的是同一件事。
%
% 取值先过一趟 \texttt{skip} 寄存器，是为了让两条路吃同一种写法：调用方按
% \hologo{LaTeX} 的胶写法给 \texttt{12pt~\cs{@plus}~0.08\cs{baselineskip}}，
% \pkg{xeCJK} 那边原样贴进 \cs{hskip}，\pkg{luatexja} 那边由寄存器吃掉再交出去。
%    \begin{macrocode}
\sys_if_engine_luatex:TF
  {
    \skip_new:N \l__hit_engine_glue_skip
    \cs_new_protected:Npn \hit@engine@cjkglue@zero:
      { \ltjsetparameter { kanjiskip = \c_zero_skip } }
    \cs_new_protected:Npn \hit@engine@cjkglue@set:nn #1#2
      {
        \skip_set:Nn \l__hit_engine_glue_skip {#1}
        \ltjsetparameter { kanjiskip = \l__hit_engine_glue_skip }
        \skip_set:Nn \l__hit_engine_glue_skip {#2}
        \ltjsetparameter { xkanjiskip = \l__hit_engine_glue_skip }
      }
  }
  {
    \cs_new_protected:Npn \hit@engine@cjkglue@zero:
      {
        \cs_set_eq:NN \CJKglue \prg_do_nothing:
        \cs_set_eq:NN \CJKecglue \prg_do_nothing:
      }
    \cs_new_protected:Npn \hit@engine@cjkglue@set:nn #1#2
      {
        \cs_set:Npx \CJKglue { \exp_not:N \hskip \exp_not:n {#1} }
        \cs_set:Npx \CJKecglue { \exp_not:N \hskip \exp_not:n {#2} }
      }
  }
%    \end{macrocode}
%
% 有几处只想按住汉字之间那一段，中西文之间的照旧，所以再给一个只管汉字间的。
%    \begin{macrocode}
\sys_if_engine_luatex:TF
  {
    \cs_new_protected:Npn \hit@engine@cjkglue@zero@cc:
      { \ltjsetparameter { kanjiskip = \c_zero_skip } }
  }
  {
    \cs_new_protected:Npn \hit@engine@cjkglue@zero@cc:
      { \cs_set_eq:NN \CJKglue \prg_do_nothing: }
  }
%    \end{macrocode}
%
% 反过来只关中西文之间那一段的也有一处：博士后目录条目里的“第1章”，数字两边不空
% （见 \file{hit-final-toc.dtx}）。\pkg{luatexja} 这边关的是自动插入，不是把
% \texttt{xkanjiskip} 设成零：页眉页脚那条路上 \pkg{docxlike} 会按字身上记的属性把
% 汉西隙的宽改回去，根本不插就改不着。
%    \begin{macrocode}
\sys_if_engine_luatex:TF
  {
    \cs_new_protected:Npn \hit@engine@ecglue@off:
      { \ltjsetparameter { autoxspacing = false } }
  }
  {
    \cs_new_protected:Npn \hit@engine@ecglue@off:
      { \cs_set_eq:NN \CJKecglue \prg_do_nothing: }
  }
%    \end{macrocode}
%
% \subsection{按汉字排一个字符}
%
% \cs{hit@engine@cjk@char:n}\marg{字符} 把一个字符按汉字排：中文字体、汉字类，只管这一处。
% 整类归汉字用上面的 \cs{hit@engine@cjk@chars:n}（\texttt{×} 就是那样一律归的）；
% 这一个给个别地方点名用。参数收字符本身，不收码位：配置层的取值在
% \pkg{expl3} 之外记号化，\texttt{"D7} 那种十六进制在那儿会碎。
%
% \hologo{XeLaTeX} 这边在\emph{分组里}把这个字符临时声明进汉字类——
% \cs{XeTeXcharclass} 赋值随分组，出组即还原，正文里的乘号一根汗毛不动
% （实测组内宋体、组外 Times）。\pkg{xeCJK} 的 \cs{CJKsymbol} 帮不上忙，
% 它就是 \cs{use:n}，字体切换在类对的记号表里。\pkg{luatexja} 有现成的
% \cs{ltjjachar}，按 JAchar 排一个码位。
%    \begin{macrocode}
\sys_if_engine_luatex:TF
  {
    \cs_new_protected:Npn \hit@engine@cjk@char:n #1
      { \ltjjachar `#1 \scan_stop: }
  }
  {
    \cs_new_protected:Npn \hit@engine@cjk@char:n #1
      {
        \group_begin:
        \xeCJKDeclareCharClass { CJK } { \int_value:w `#1 }
        #1
        \group_end:
      }
  }
\cs_new_eq:NN \hit@engine@cjk@char \hit@engine@cjk@char:n
%    \end{macrocode}
%
% \subsection{今天（东八区）}
%
% \cs{g\_\_hit\_engine\_today\_year\_int}、\cs{\ldots month\ldots}、
% \cs{\ldots day\ldots} 是东八区的今天。\pkg{expl3} 的 \cs{c\_sys\_year\_int}
% 一族是编译机器的本地时间，不带时区；引擎各有一条拿得到时区的路：
%
% \hologo{XeLaTeX} 的 \cs{creationdate} 给
% \texttt{D:20260830001418+08'00'} 这种串，本地时刻连着时区偏移。按偏移折回
% UTC 再加八小时，跨天最多挪一天（偏移在 $-12$ 到 $+14$ 小时之间），月初年初
% 的进退位按月长算，闰年按四百年规则。串的末尾也可能是 \texttt{Z}（UTC）或者
% 什么都没有，都当零偏移。\texttt{SOURCE\_DATE\_EPOCH} 一样管用——
% \texttt{FORCE\_SOURCE\_DATE} 下 \cs{creationdate} 从环境变量来，回归测试
% 钉日期的做法不受影响。
%
% \hologo{LuaLaTeX} 没有 \cs{creationdate}，一句
% \texttt{os.date} 拿 UTC 加八小时，不用自己算进退位。格式串里不能出现
% \texttt{\%}——它在生成的类文件里是注释符，行尾会被整个吃掉、花括号失衡，
% \cs{sys\_if\_engine\_luatex:TF} 扫参数时把后面整段定义吞进这个分支（踩过：
% 三个解析函数全成了残废，东八区悄悄退回本地日期）。所以取
% \texttt{os.date("!*t")} 的表字段，一个 \texttt{\%} 都不用。
%
% 两条路都不通（\cs{creationdate} 不在或者串认不出）就退回本地时间。
%    \begin{macrocode}
\int_new:N \g__hit_engine_today_year_int
\int_new:N \g__hit_engine_today_month_int
\int_new:N \g__hit_engine_today_day_int
\int_gset_eq:NN \g__hit_engine_today_year_int \c_sys_year_int
\int_gset_eq:NN \g__hit_engine_today_month_int \c_sys_month_int
\int_gset_eq:NN \g__hit_engine_today_day_int \c_sys_day_int
%    \end{macrocode}
%
% 解析放在用不用得上都无害的形状里：串长不够、打头不是 \texttt{D:}，
% 就什么都不改。
%    \begin{macrocode}
\cs_new_protected:Npn \__hit_engine_today_parse:n #1
  { \exp_args:Nx \__hit_engine_today_split:n { \tl_to_str:n {#1} } }
\int_new:N \l__hit_engine_offset_int
\int_new:N \l__hit_engine_total_int
\cs_new_protected:Npn \__hit_engine_today_split:n #1
  {
    \int_compare:nNnT { \str_count:n {#1} } > { 15 }
      {
        \str_if_eq:eeT { \str_range:nnn {#1} { 1 } { 2 } } { D: }
          {
            \__hit_engine_today_set:xxxxxx
              { \str_range:nnn {#1} {  3 } {  6 } }
              { \str_range:nnn {#1} {  7 } {  8 } }
              { \str_range:nnn {#1} {  9 } { 10 } }
              { \str_range:nnn {#1} { 11 } { 12 } }
              { \str_range:nnn {#1} { 13 } { 14 } }
              { \str_range:nnn {#1} { 17 } { \str_count:n {#1} } }
          }
      }
  }
%    \end{macrocode}
%
% 偏移那一截是 \texttt{+08'00'}、\texttt{-06'30'}、\texttt{Z} 或空。
% 折成分钟，认不出当零。
%    \begin{macrocode}
\cs_new_protected:Npn \__hit_engine_today_offset:n #1
  {
    \int_zero:N \l__hit_engine_offset_int
    \int_compare:nNnT { \str_count:n {#1} } > { 5 }
      {
        \bool_lazy_or:nnT
          { \str_if_eq_p:ee { \str_head:n {#1} } { + } }
          { \str_if_eq_p:ee { \str_head:n {#1} } { - } }
          {
            \int_set:Nn \l__hit_engine_offset_int
              {
                \str_range:nnn {#1} { 2 } { 3 } * 60
                + \str_range:nnn {#1} { 5 } { 6 }
              }
            \str_if_eq:eeT { \str_head:n {#1} } { - }
              {
                \int_set:Nn \l__hit_engine_offset_int
                  { - \l__hit_engine_offset_int }
              }
          }
      }
  }
%    \end{macrocode}
%
% 本地时刻减本地偏移是 UTC，再加 480 分钟是东八区。整除 1440 得跨天数，
% 只会是 $-1$、$0$ 或 $+1$；日在月里进退位，月长按闰年查。
%    \begin{macrocode}
\cs_new_protected:Npn \__hit_engine_today_set:nnnnnn #1#2#3#4#5#6
  {
    \__hit_engine_today_offset:n {#6}
    \int_gset:Nn \g__hit_engine_today_year_int {#1}
    \int_gset:Nn \g__hit_engine_today_month_int {#2}
    \int_gset:Nn \g__hit_engine_today_day_int {#3}
    \int_set:Nn \l__hit_engine_total_int
      { #4 * 60 + #5 - \l__hit_engine_offset_int + 480 }
    \int_compare:nNnT { \l__hit_engine_total_int } < { 0 }
      { \__hit_engine_today_shift:n { -1 } }
    \int_compare:nNnT { \l__hit_engine_total_int } > { 1439 }
      { \__hit_engine_today_shift:n { 1 } }
  }
\cs_generate_variant:Nn \__hit_engine_today_set:nnnnnn { xxxxxx }
\cs_new:Npn \__hit_engine_month_days:nn #1#2
  {
    \int_case:nnF {#2}
      {
        { 2 }
          {
            \int_compare:nNnTF
              { \int_mod:nn {#1} { 4 } +
                \int_compare_p:nNn { \int_mod:nn {#1} { 100 } } = { 0 }
                - \int_compare_p:nNn { \int_mod:nn {#1} { 400 } } = { 0 } }
              = { 0 }
              { 29 } { 28 }
          }
        { 4 } { 30 } { 6 } { 30 } { 9 } { 30 } { 11 } { 30 }
      }
      { 31 }
  }
\cs_new_protected:Npn \__hit_engine_today_shift:n #1
  {
    \int_gadd:Nn \g__hit_engine_today_day_int {#1}
    \int_compare:nNnT { \g__hit_engine_today_day_int } < { 1 }
      {
        \int_gadd:Nn \g__hit_engine_today_month_int { -1 }
        \int_compare:nNnT { \g__hit_engine_today_month_int } < { 1 }
          {
            \int_gset:Nn \g__hit_engine_today_month_int { 12 }
            \int_gadd:Nn \g__hit_engine_today_year_int { -1 }
          }
        \int_gset:Nn \g__hit_engine_today_day_int
          {
            \__hit_engine_month_days:nn
              { \g__hit_engine_today_year_int }
              { \g__hit_engine_today_month_int }
          }
      }
    \int_compare:nNnT
      { \g__hit_engine_today_day_int } >
      {
        \__hit_engine_month_days:nn
          { \g__hit_engine_today_year_int }
          { \g__hit_engine_today_month_int }
      }
      {
        \int_gset:Nn \g__hit_engine_today_day_int { 1 }
        \int_gadd:Nn \g__hit_engine_today_month_int { 1 }
        \int_compare:nNnT { \g__hit_engine_today_month_int } > { 12 }
          {
            \int_gset:Nn \g__hit_engine_today_month_int { 1 }
            \int_gadd:Nn \g__hit_engine_today_year_int { 1 }
          }
      }
  }
%    \end{macrocode}
%
% 算一次就定死。必须放在定义之后：写在前面时加载期的 Undefined control sequence
% 被 nonstopmode 咽掉，东八区悄悄退回本地日期，只有探针测试抓得出。
%    \begin{macrocode}
\sys_if_engine_luatex:TF
  {
    \int_gset:Nn \g__hit_engine_today_year_int
      { \directlua { tex.sprint ( os.date ( "!*t" , os.time ( ) + 28800 ) .year ) } }
    \int_gset:Nn \g__hit_engine_today_month_int
      { \directlua { tex.sprint ( os.date ( "!*t" , os.time ( ) + 28800 ) .month ) } }
    \int_gset:Nn \g__hit_engine_today_day_int
      { \directlua { tex.sprint ( os.date ( "!*t" , os.time ( ) + 28800 ) .day ) } }
  }
  {
    \cs_if_exist:NT \creationdate
      { \exp_args:Nx \__hit_engine_today_parse:n { \creationdate } }
    \sys_if_engine_pdftex:T
      { \exp_args:Nx \__hit_engine_today_parse:n { \pdfcreationdate } }
  }
%    \end{macrocode}
%
% \subsection{一串文字里有没有汉字（\hologo{pdfTeX}）}
%
% \cs{hit@engine@pdftex@cjk@scan:nN}\marg{文字}\meta{布尔}：文字里有汉字就把布尔设真。两个 Unicode
% 引擎上汉字是一个字符记号，调用方自己逐记号取码位比 \texttt{U+2E7F}；\hologo{pdfTeX} 下
% 汉字是几个字节，展开后又成了 \cs{CTEX@char@nnn} 那样的宏，那一套取不到码位。这里先用
% \cs{text\_purify:n} 还原成裸字节（\pkg{ctex} 给那几个宏登记了还原办法），再看有没有
% 0xE3 起头的三字节字（\texttt{U+3000} 以上：汉字、假名、全角标点）。
%    \begin{macrocode}
\str_new:N \l__hit_engine_scan_str
\cs_new_protected:Npn \hit@engine@pdftex@cjk@scan:nN #1#2
  {
    \str_set:Ne \l__hit_engine_scan_str { \text_purify:n {#1} }
    \str_map_inline:Nn \l__hit_engine_scan_str
      {
        \int_compare:nNnT { `##1 } > { 226 }
          { \bool_set_true:N #2 \str_map_break: }
      }
  }
%    \end{macrocode}
%
% \subsection{按字拆开（\hologo{pdfTeX}）}
%
% \cs{hit@engine@pdftex@chars:nN}\marg{文字}\meta{序列} 把文字按字拆成一串，每项一个字。
% 两个 Unicode 引擎上一个汉字就是一个记号，\cs{tl\_map\_inline:nn} 逐项走就是逐字走；
% \hologo{pdfTeX} 下 \pkg{CJKutf8} 把 UTF-8 的每个字节做成一个活动字符，一个汉字三个记号，
% 逐记号走会把一个字劈开（首字节展开时去抓后面两个字节，抓到的是映射本身的记号）。这里先
% 用 \cs{text\_purify:n} 把文字还原成裸字节，再照当前的 catcode 重新记号化（字节回到活动字符），
% 然后把 0x80--0xBF 的后续字节并进前一个字。\cs{text\_purify:n} 顺手把命令也剥掉了，
% 所以只拿它拆字，不拿它排带格式的文字。
%    \begin{macrocode}
\tl_new:N \l__hit_engine_chars_tl
\tl_new:N \l__hit_engine_char_tl
\cs_new_protected:Npn \hit@engine@pdftex@chars:nN #1#2
  {
    \seq_clear:N #2
    \tl_clear:N \l__hit_engine_char_tl
    \str_set:Ne \l__hit_engine_scan_str { \text_purify:n {#1} }
    \tl_set_rescan:NnV \l__hit_engine_chars_tl { } \l__hit_engine_scan_str
    \tl_map_inline:Nn \l__hit_engine_chars_tl
      {
        \bool_lazy_and:nnTF
          { \tl_if_single_token_p:n {##1} }
          { \token_if_active_p:N ##1 }
          {
            \int_compare:nNnTF { `##1 } < { "C0 }
              { \tl_put_right:Nn \l__hit_engine_char_tl {##1} }
              {
                \tl_if_empty:NF \l__hit_engine_char_tl
                  { \seq_put_right:NV #2 \l__hit_engine_char_tl }
                \tl_set:Nn \l__hit_engine_char_tl {##1}
              }
          }
          {
            \tl_if_empty:NF \l__hit_engine_char_tl
              { \seq_put_right:NV #2 \l__hit_engine_char_tl }
            \tl_set:Nn \l__hit_engine_char_tl {##1}
          }
      }
    \tl_if_empty:NF \l__hit_engine_char_tl
      { \seq_put_right:NV #2 \l__hit_engine_char_tl }
  }
\cs_generate_variant:Nn \hit@engine@pdftex@chars:nN { V }
%    \end{macrocode}
%
% \cs{hit@engine@group@chars:N}\meta{记号列} 把记号列里的汉字各包成一组，别的记号不动，之后逐项
% 走（\cs{tl\_map\_inline:Nn}、\cs{tl\_item:Nn}）就是逐字走。两个 Unicode 引擎上一个汉字本来就是
% 一个记号，什么都不做。\hologo{pdfTeX} 下汉字有两种样子：还没展开的是一串活动字节（首字节
% 0xC0 以上，后面跟 0x80--0xBF），\cs{tl\_set:Ne} 展开过的是 \pkg{ctex} 的
% \cs{CTEX@char@nnn} 带三个十进制字节参数（两字节、四字节的同理）。与上面那个不同，这里不
% purify，夹在字间的命令（\cs{hit@sp} 之类）都留着。
%
% \cs{hit@engine@char@key:nN}\marg{一个字}\meta{记号列} 取查表用的键：属性表的键是字符串，
% \hologo{pdfTeX} 下写在文件里的字存成裸字节，展开过的那种样子要先 purify 回去才对得上。
%    \begin{macrocode}
\sys_if_engine_pdftex:TF
  {
    \regex_const:Nn \c__hit_engine_char_regex
      {
        \c{CTEX@char@nnnn} (\cB. \d+ \cE.){4}
        | \c{CTEX@char@nnn} (\cB. \d+ \cE.){3}
        | \c{CTEX@char@nn} (\cB. \d+ \cE.){2}
        | \cA[\x{C0}-\x{FF}] \cA[\x{80}-\x{BF}]*
      }
    \cs_new_protected:Npn \hit@engine@group@chars:N #1
      { \regex_replace_all:NnN \c__hit_engine_char_regex { \cB\{ \0 \cE\} } #1 }
    \cs_new_protected:Npn \hit@engine@char@key:nN #1#2
      { \tl_set:Ne #2 { \text_purify:n {#1} } }
  }
  {
    \cs_new_protected:Npn \hit@engine@group@chars:N #1 { }
    \cs_new_protected:Npn \hit@engine@char@key:nN #1#2 { \tl_set:Nn #2 {#1} }
  }
%    \end{macrocode}
%
% \subsection{定宽分散对齐}
%
% \cs{hit@engine@spread@to:nn}\marg{宽度}\marg{文本} 把文本分散对齐到给定宽度，
% 撑开的是汉字之间那段胶——所以 \texttt{学院（部} 不会在左括号后面多出一格。
% 对外的名字是 \file{src/engine/hit-spread.dtx} 的 \cs{hit@spread@to}。
%
% \pkg{xeCJK} 那边现成：\pkg{xeCJKfntef} 的 \env{CJKfilltwosides}。
% \pkg{luatexja} 那边没有对应的环境，照 \cls{thuthesis} 的写法自己搭：
% \texttt{kanjiskip} 本来就是那段胶，把它的伸展量设足再套进定宽
% \cs{makebox}，撑开的位置与量都一样。
%
% \textbf{伸展量取 \texttt{filll}，盒子取 \oarg{l} 而不是 \oarg{s}。} 这两条是配套的：
% \cs{makebox}\oarg{宽度}\oarg{l} 会在内容右边垫一个 \cs{hfil}，而 \texttt{filll}
% 比 \texttt{fil} 高一阶，字距那段胶总是先被撑开、垫的那个吃不到伸展量。好处是
% 内容里一段汉字字距都没有时（纯西文标签）不会报 underfull，而是老老实实左对齐。
% \cls{thuthesis} 的 \cs{thu@fixed@box} 就是这么写的，取值也照抄它的
% \texttt{0pt~plus~2filll~minus~1filll}。
%
% \textbf{内容撑不满才有意义。} 两条路都不处理内容比给定宽度还宽的情形：
% \env{CJKfilltwosides} 是个 \env{minipage}，会折行；这边的 \cs{makebox} 会溢出。
% 要「压进这个宽度」写 \cs{makebox}\oarg{宽度}\oarg{s}，见 \file{hit-spread.dtx}。
%
% ^^A TODO(v3.2b)：\cls{thuthesis} 两个引擎共用这一套，不碰 \env{CJKfilltwosides}。
% ^^A 这边 \hologo{XeLaTeX} 仍走原环境，是为了不动既有的排版基线。等 make smoke
% ^^A 验过两者逐页相同，再把 \hologo{XeLaTeX} 也并过来，这个分支就能收掉。
%    \begin{macrocode}
\sys_if_engine_luatex:TF
  {
    \cs_new_protected:Npn \hit@engine@spread@to:nn #1#2
      {
        \group_begin:
          \ltjsetparameter { kanjiskip = { 0pt~plus~2filll~minus~1filll } }
          \makebox [ #1 ] [ l ] {#2}
        \group_end:
      }
  }
  {
    \cs_new_protected:Npn \hit@engine@spread@to:nn #1#2
      { \begin { CJKfilltwosides } [b] {#1} #2 \end { CJKfilltwosides } }
  }
%    \end{macrocode}
%
% \hologo{pdfTeX} 下 \pkg{CJK} 在汉字之间插的也是 \cs{CJKglue}，照 \hologo{LuaLaTeX} 那一支的办法撑它。
%    \begin{macrocode}
\sys_if_engine_pdftex:T
  {
    \cs_set_protected:Npn \hit@engine@spread@to:nn #1#2
      {
        \group_begin:
          \cs_set:Npn \CJKglue { \hskip 0pt~plus~2filll~minus~1filll \scan_stop: }
          \makebox [ #1 ] [ l ] {#2}
        \group_end:
      }
  }
%    \end{macrocode}
%
% \subsection{当前中文字体里有没有这个字}
%
% \cs{hit@engine@if@cjk@glyph:nTF}\marg{码位} 问当前中文字体有没有这个码位的
% 字形。\cs{quan} 拿它决定是用字体自带的带圈数字，还是自己画一个。
%
% \pkg{xeCJK} 那边 \cs{XeTeXcharglyph} 问的是「当前字体」，而
% \cs{\_\_hit\_use\_cjk\_font:} 已经把当前字体切成了中文字体，直接问即可。
%
% \pkg{luatexja} 那边中文字体不是「当前字体」——它另挂一个和文字体，字体号存在
% \cs{ltj@curjfnt} 属性里，\cs{fontname}\cs{font} 拿到的仍是西文字体。所以得绕到
% Lua 那边按属性取出字体号再查字形表。属性没设（不在中文字体环境里）时当作没有。
%    \begin{macrocode}
\sys_if_engine_luatex:TF
  {
    \cs_new:Npn \__hit_engine_glyph:n #1
      {
        \tex_directlua:D
          {
            local~a = luatexbase.attributes['ltj@curjfnt']
            local~i = a~and~tex.attribute[a]
            local~f = (i~and~i >= 0)~and~font.getfont(i)
            if~f~and~f.characters~and~f.characters[\int_eval:n {#1}]
              then~tex.print('1')~else~tex.print('0')~end
          }
      }
    \prg_new_conditional:Npnn \hit@engine@if@cjk@glyph:n #1 { p , TF }
      {
        \int_compare:nNnTF { \__hit_engine_glyph:n {#1} } > { 0 }
          { \prg_return_true: } { \prg_return_false: }
      }
  }
  {
    \prg_new_conditional:Npnn \hit@engine@if@cjk@glyph:n #1 { p , TF }
      {
        \int_compare:nNnTF { \XeTeXcharglyph \int_eval:n {#1} } > { 0 }
          { \prg_return_true: } { \prg_return_false: }
      }
  }
%    \end{macrocode}
%
% \subsection{把当前字体切成中文字体}
%
% \cs{hit@engine@use@cjk@font:} 让接下来的西文字符也从中文字体取字形。
% \cs{quan} 画圈时圈里的数字要跟正文宋体一致，得真的切过去，光选中文字族不够。
%
% \pkg{xeCJK} 底下中文字族就是个 NFSS 字族，\cs{usefont} 切过去即可。
%
% \pkg{luatexja} 底下办不到，那边是空操作。和文字体不走 NFSS 字族这条路，
% 字体号另挂在 \cs{ltj@curjfnt} 属性上，拿 \cs{CJK@family} 去 \cs{usefont}
% 只会落回默认西文字体（实测 \cs{fontname}\cs{font} 仍是 \pkg{lmroman}）。
% 把数字那段码位划归和文倒是能达到目的，但 \cs{ltjdefcharrange} 是全局的，
% 写在组里也退不回来，整篇的数字都会跟着换字体，代价太大。
%
% 后果只落在 \cs{quan} 画圈那一支上：圈里的数字取自西文字体，比
% \hologo{XeLaTeX} 那边细一点。字体自带带圈字符时走的是另一支，不受影响，
% 而常见的几档中文字体都自带 \texttt{①}--\texttt{⑳}，所以多数文档看不出差别。
%    \begin{macrocode}
\sys_if_engine_luatex:TF
  { \cs_new_protected:Npn \hit@engine@use@cjk@font: { } }
  {
    \cs_new_protected:Npn \hit@engine@use@cjk@font:
      { \usefont { \f@encoding } { \CJK@family } { m } { n } }
  }
%    \end{macrocode}
%
% \hologo{pdfTeX} 下中文字族是 \pkg{CJK} 的 C70 编码，拿它配西文编码是个不存在的字体形状，
% 切不过去，这里什么都不做；问字形也问不了，带圈数字只认中易宋体一定有的 ①--⑳，文鼎那一档
% 一律画。\cs{hit@engine@char:n}\marg{码位} 排一个字：两个 Unicode 引擎是 \cs{char}，\hologo{pdfTeX}
% 下把码位拆成 UTF-8 的三个字节交给 \pkg{ctex} 的 \cs{CTEX@char@nnn}（活动字符展开后就是它）。
%    \begin{macrocode}
\int_new:N \l__hit_engine_cp_int
\sys_if_engine_pdftex:TF
  {
    \cs_set_protected:Npn \hit@engine@use@cjk@font: { }
    \prg_set_conditional:Npnn \hit@engine@if@cjk@glyph:n #1 { p , TF }
      {
        \bool_lazy_and:nnTF
          { ! \bool_if_p:N \g__hit_pdftex_arphic_bool }
          {
            \bool_lazy_and_p:nn
              { ! \int_compare_p:nNn {#1} < { "2460 } }
              { ! \int_compare_p:nNn {#1} > { "2473 } }
          }
          { \prg_return_true: } { \prg_return_false: }
      }
    \cs_new_protected:Npn \hit@engine@char:n #1
      {
        \int_set:Nn \l__hit_engine_cp_int {#1}
        \use:e
          {
            \exp_not:N \CTEX@char@nnn
              { \int_eval:n { 224 + \int_div_truncate:nn { \l__hit_engine_cp_int } { 4096 } } }
              {
                \int_eval:n
                  { 128 + \int_mod:nn { \int_div_truncate:nn { \l__hit_engine_cp_int } { 64 } } { 64 } }
              }
              { \int_eval:n { 128 + \int_mod:nn { \l__hit_engine_cp_int } { 64 } } }
          }
      }
  }
  {
    \cs_new_protected:Npn \hit@engine@char:n #1
      { \tex_char:D \int_eval:n {#1} \scan_stop: }
  }
%    \end{macrocode}
%
% \subsection{中西文之间的间距（\hologo{pdfTeX}）}
%
% Word 在汉字与西文字母、数字之间自动留一段空（\texttt{autoSpaceDE}/\texttt{DN}），
% \pkg{docxlike} 把它算好放进 \cs{CJKecglue}。\pkg{xeCJK} 与 \pkg{luatexja} 按字符类自己在
% 交界处插这段胶，\pkg{CJK} 宏包没有这个机制，只在前一个也是汉字时插 \cs{CJKglue}。这里照
% \pkg{gxuthesis} 的办法补上，包的是 \pkg{ctex} 排三字节、四字节字的 \cs{CTEX@char@nnn} 与
% \cs{CTEX@char@nnnn}（汉字的活动字节展开后就是它们）：
% \begin{itemize}
% \item 左边界：排字之前看上一个节点。是字符、连字（上一个是西文字）、行内公式，或者不是
%   \pkg{CJK} 记号的 kern（\cs{textit} 之后的斜体校正）就插胶。前面是字符的再看
%   \cs{spacefactor}：只有字母与数字（999、1000）才插。kern 也按前一个字符看，直立体的斜体
%   校正是一个零宽的 kern，1--3\,sp 的是记号。\texttt{( [ ) ] / -} 的 sfcode 改成
%   998，\cs{frenchspacing} 把句号那几个拨回 1000 时也跟着改，Word 在这些符号旁不留空。
%   取小于 1000 的数是因为大于 1000 的跟在大写字母后面会被 \hologo{TeX} 压回 1000；
%   998 对词间胶的影响是千分之二。
% \item 右边界：排完之后吞掉空格（与 \pkg{xeCJK} 的 \texttt{CJKspace=false} 一致），看下一个
%   记号。字母、数字、\verb|$| 插胶；控制序列里几个排字体与引用的，先插胶再垫一个 3\,sp 的
%   kern 当记号，里面排出来的若是汉字，它的左边界认出记号把胶撤掉。
%   别的控制序列能展开、又不是 \cs{protected} 的，先展开一层再看（\verb|\the\month| 展开出数字，
%   日期里“年”后面就有胶），最多十六层；展开出下一个汉字（活动字节展开头一步是 \pkg{CJK} 的
%   \cs{CJK@@@}）或者碰到不能展开的就停，不插。\cs{label}、\cs{index} 这类先落一个 whatsit，
%   插了就撤不掉，所以不走试探那一条。先展开一层与 \hologo{TeX} 自己往下走时展开的是同一个记号，
%   次序不变。
% \item 全角标点两边都不插。码位落在通用标点、中日韩符号与标点、全角 ASCII 标点这几段的
%   是标点；\pkg{CJKpunct} 自己接管的弯引号、破折号、省略号不走 \cs{CTEX@char@nnn}，本来就不插。
% \item 行内代码与 \env{verbatim} 里不插，与 \pkg{xeCJK} 一样。
% \end{itemize}
% 没有 \pkg{docxlike} 设过的地方（封面那一类整页）\cs{CJKecglue} 取 0.25\,em，与
% \pkg{luatexja} 的缺省相当。
%
% \verb|~| 也在这里还原成 \cs{nobreakspace}。\pkg{ctex} 在导言区末尾调 \pkg{CJK} 的
% \cs{CJKtilde}，把它换成一个会看两边是不是汉字的宏，不健壮：英文题注前缀 \texttt{Fig.\textasciitilde}
% 在 \pkg{ccaption} 的 \cs{edef} 里一展开就报错。另两个引擎下它本来就是 \cs{nobreakspace}。
%    \begin{macrocode}
\bool_new:N \l_hit_engine_ec_off_bool
\sys_if_engine_pdftex:T
  {
    \bool_new:N \l__hit_engine_ec_punct_bool
    \int_new:N \l__hit_engine_ec_cp_int
    \tl_const:Nn \c__hit_engine_ec_start_tl
      {
        \textbf \textit \textsl \textrm \textsf \texttt \textup \textmd
        \textnormal \textsc \emph \ref \eqref \pageref \autoref \num \SI
        \si \qty \unit \LaTeX \TeX \LaTeXe
      }
    \cs_new_protected:Npn \__hit_engine_ec_glue:
      {
        \cs_if_exist:NTF \CJKecglue
          { \CJKecglue }
          { \tex_hskip:D 0.25em plus 0.125em minus 0.08em \scan_stop: }
      }
    \cs_new_protected:Npn \__hit_engine_ec_punct:n #1
      {
        \int_set:Nn \l__hit_engine_ec_cp_int {#1}
        \bool_set:Nn \l__hit_engine_ec_punct_bool
          {
            \int_compare_p:n { "2010 <= \l__hit_engine_ec_cp_int <= "205F }
            || \int_compare_p:n { "3000 <= \l__hit_engine_ec_cp_int <= "303F }
            || \int_compare_p:n { "FE10 <= \l__hit_engine_ec_cp_int <= "FE6F }
            || \int_compare_p:n { "FF01 <= \l__hit_engine_ec_cp_int <= "FF0F }
            || \int_compare_p:n { "FF1A <= \l__hit_engine_ec_cp_int <= "FF20 }
            || \int_compare_p:n { "FF3B <= \l__hit_engine_ec_cp_int <= "FF40 }
            || \int_compare_p:n { "FF5B <= \l__hit_engine_ec_cp_int <= "FF65 }
          }
      }
    \cs_new_protected:Npn \__hit_engine_ec_nnn:nnn #1#2#3
      {
        \__hit_engine_ec_punct:n
          { (#1 - "E0) * "1000 + (#2 - "80) * "40 + #3 - "80 }
        \__hit_engine_ec_wrap:n { \__hit_engine_ec_orig_nnn:nnn {#1} {#2} {#3} }
      }
    \cs_new_protected:Npn \__hit_engine_ec_nnnn:nnnn #1#2#3#4
      {
        \bool_set_false:N \l__hit_engine_ec_punct_bool
        \__hit_engine_ec_wrap:n
          { \__hit_engine_ec_orig_nnnn:nnnn {#1} {#2} {#3} {#4} }
      }
    \cs_new_protected:Npn \__hit_engine_ec_wrap:n #1
      {
        \mode_if_math:TF {#1}
          {
            \mode_leave_vertical:
            \__hit_engine_ec_left:
            #1
            \bool_if:NF \l_hit_engine_ec_off_bool
              {
                \int_zero:N \l__hit_engine_ec_expand_int
                \__hit_engine_ec_right:
              }
          }
      }
%    \end{macrocode}
%
% 左边界。先撤右边界垫过的胶：上一个节点是 3\,sp 的记号就连它前面那段胶一起拿掉。
%    \begin{macrocode}
    \cs_new_protected:Npn \__hit_engine_ec_left:
      {
        \bool_lazy_and:nnT
          { \int_compare_p:nNn { \tex_lastnodetype:D } = { 12 } }
          { \int_compare_p:nNn { \tex_lastkern:D } = { 3 } }
          { \tex_unkern:D \tex_unskip:D }
        \bool_lazy_or:nnF
          { \l__hit_engine_ec_punct_bool }
          { \l_hit_engine_ec_off_bool }
          {
            \int_case:nn { \tex_lastnodetype:D }
              {
                { 0 } { \__hit_engine_ec_after_char: }
                { 7 } { \__hit_engine_ec_after_char: }
                { 10 } { \__hit_engine_ec_glue: }
                { 12 }
                  {
                    \bool_lazy_and:nnF
                      { \int_compare_p:nNn { \tex_lastkern:D } > { 0 } }
                      { \int_compare_p:nNn { \tex_lastkern:D } < { 4 } }
                      { \__hit_engine_ec_after_char: }
                  }
              }
          }
      }
    \cs_new_protected:Npn \__hit_engine_ec_after_char:
      {
        \int_compare:nNnT { \tex_spacefactor:D } > { 998 }
          {
            \int_compare:nNnT { \tex_spacefactor:D } < { 1001 }
              { \__hit_engine_ec_glue: }
          }
      }
%    \end{macrocode}
%
% 右边界。下一个记号只比含义，不取下来：前瞻碰到单元格的 \verb|&| 时 \hologo{TeX} 已经把
% 模板的后半截插了进来，取下来再放回去就把单元格排乱了。空格自己吞，不用
% \cs{peek\_remove\_spaces:n}：它在对齐安全组里前瞻，出组后再前瞻一次，单元格末尾的
% \verb|&| 就报“Extra alignment tab”，纯 \cls{article} 里也是这样。
%    \begin{macrocode}
    \cs_new_protected:Npn \__hit_engine_ec_right:
      { \peek_after:Nw \__hit_engine_ec_skip: }
    \cs_new_protected:Npn \__hit_engine_ec_skip:
      {
        \token_if_eq_meaning:NNTF \l_peek_token \c_space_token
          {
            \tex_afterassignment:D \__hit_engine_ec_right:
            \cs_set_eq:NN \__hit_engine_ec_tmp:w
          }
          { \__hit_engine_ec_peek: }
      }
    \int_new:N \l__hit_engine_ec_expand_int
    \cs_new_protected:Npn \__hit_engine_ec_expand:
      {
        \int_incr:N \l__hit_engine_ec_expand_int
        \exp_after:wN \__hit_engine_ec_right:
      }
    \cs_new_protected:Npn \__hit_engine_ec_peek:
      {
        \cs_set_eq:NN \__hit_engine_ec_next: \prg_do_nothing:
        \bool_if:NF \l__hit_engine_ec_punct_bool
          {
            \bool_lazy_or:nnTF
              { \token_if_letter_p:N \l_peek_token }
              { \token_if_math_toggle_p:N \l_peek_token }
              { \__hit_engine_ec_glue: }
              {
                \token_if_other:NTF \l_peek_token
                  {
                    \tl_map_inline:nn { 0123456789 }
                      {
                        \token_if_eq_charcode:NNT \l_peek_token ##1
                          { \__hit_engine_ec_glue: \tl_map_break: }
                      }
                  }
                  { \token_if_cs:NT \l_peek_token { \__hit_engine_ec_cs: } }
              }
          }
        \__hit_engine_ec_next:
      }
    \bool_new:N \l__hit_engine_ec_start_bool
    \cs_new_protected:Npn \__hit_engine_ec_cs:
      {
        \bool_set_false:N \l__hit_engine_ec_start_bool
        \tl_map_inline:Nn \c__hit_engine_ec_start_tl
          {
            \token_if_eq_meaning:NNT \l_peek_token ##1
              {
                \bool_set_true:N \l__hit_engine_ec_start_bool
                \tl_map_break:
              }
          }
        \bool_if:NT \l__hit_engine_ec_start_bool
          {
            \__hit_engine_ec_glue:
            \tex_kern:D 3 sp \scan_stop:
          }
        \bool_lazy_all:nT
          {
            { ! \l__hit_engine_ec_start_bool }
            { ! \token_if_eq_meaning_p:NN \l_peek_token \CJK@@@ }
            { \token_if_expandable_p:N \l_peek_token }
            { ! \token_if_protected_macro_p:N \l_peek_token }
            { ! \token_if_protected_long_macro_p:N \l_peek_token }
            { \int_compare_p:nNn { \l__hit_engine_ec_expand_int } < { 16 } }
          }
          { \cs_set_eq:NN \__hit_engine_ec_next: \__hit_engine_ec_expand: }
      }
%    \end{macrocode}
%
% 装上。\pkg{ctex} 在导言区末尾才开 \pkg{CJK} 环境，这里等到 \texttt{begindocument}。
%    \begin{macrocode}
    \cs_new_protected:Npn \__hit_engine_ec_sfcodes:
      {
        \clist_map_inline:nn { `( , `[ , `) , `] , `/ , `- , `. , `? , `! , `: , `; , `, }
          { \tex_sfcode:D ##1 = 998 \scan_stop: }
      }
    \cs_new_protected:Npn \__hit_engine_ec_install:
      {
        \char_set_active_eq:nN { `\~ } \nobreakspace
        \cs_set_eq:NN \__hit_engine_ec_orig_nnn:nnn \CTEX@char@nnn
        \cs_set_eq:NN \__hit_engine_ec_orig_nnnn:nnnn \CTEX@char@nnnn
        \cs_set_eq:NN \CTEX@char@nnn \__hit_engine_ec_nnn:nnn
        \cs_set_eq:NN \CTEX@char@nnnn \__hit_engine_ec_nnnn:nnnn
        \tl_put_right:Nn \verbatim@font { \bool_set_true:N \l_hit_engine_ec_off_bool }
        \tl_put_right:Nn \frenchspacing { \__hit_engine_ec_sfcodes: }
        \int_compare:nNnTF { \tex_sfcode:D `\. } = { 1000 }
          { \__hit_engine_ec_sfcodes: }
          {
            \clist_map_inline:nn { `( , `[ , `) , `] , `/ , `- }
              { \tex_sfcode:D ##1 = 998 \scan_stop: }
          }
      }
    \AddToHook { begindocument } { \__hit_engine_ec_install: }
  }
%    \end{macrocode}
%
% \subsection{装载基础类之后要补的两件事}
%
% \cs{hit@engine@postload:} 由两个 \file{*-deps.dtx} 调用。本文件整体排在
% \cs{LoadClass} 之前，那时 \pkg{luatexja} 还没装，\cs{ltjdefcharrange} 之类
% 一个都不存在，所以凡是要\emph{执行}的都推到这里，本文件里只管\emph{定义}。
% \hologo{XeLaTeX} 下它是空的。
%
% \textbf{一、带圈数字归成汉字。} \cs{quan} 排脚注号时先问字体有没有现成的
% \texttt{①}，有就直接排那个码位。\pkg{luatexja} 按区块把
% \texttt{U+2460}--\texttt{U+24FF} 归了西文，于是那一句排出来去找 Times New
% Roman，字体里没有就整个排丢了（日志里一句 Missing character，纸面上什么都没有）。
% 归成汉字之后从中文字体取，与 \hologo{XeLaTeX} 排的是同一个字形。
% \cls{thuthesis} 的 \cs{ltjdefcharrange}\marg{6} 里也把这一段划了回去。
%
% 三段码位与 \cs{\_\_hit\_quan\_slot:n} 的三段一一对应：1--20 在
% \texttt{U+2460}，21--35 在 \texttt{U+3251}，36--50 在 \texttt{U+32B1}。
%
% \hologo{XeLaTeX} 不需要这一句：那边 \cs{\_\_hit\_use\_cjk\_font:} 已经把当前
% 字体整个切成了中文字体，码位归哪一类都取得到字形。
%
% \textbf{二、下划线那一族。} 拿 \pkg{ulem} 顶出 \pkg{xeCJKfntef} 的几个命令。
%
% 类自己一处都不调这几个命令，它们是给用户写正文用的，手册里也写着。老论文里
% 写了 \cs{CJKunderline} 的，换个引擎不该编不过，所以顶一下。
%
% \textbf{顶不出来的两个。} \cs{CJKunderanyline} 与 \cs{CJKunderanysymbol} 要
% 自定义线型与符号，\pkg{ulem} 没有对得上的东西，不给。用到它们的文档在
% \hologo{LuaLaTeX} 下会报未定义，得改用 \hologo{XeLaTeX} 编。
%
% \textbf{断行行为不一样。} \pkg{xeCJKfntef} 认汉字，\pkg{ulem} 是按盒子拆的，
% 长段落里的断点位置会与 \hologo{XeLaTeX} 有出入。短语上看不出来，整段加下划线
% 就明显了。
%    \begin{macrocode}
\sys_if_engine_luatex:TF
  {
    \cs_new_protected:Npn \hit@engine@postload:
      {
        \hit@engine@cjk@chars:n
          { "2460 - "2473 , "3251 - "325F , "32B1 - "32BF }
        \cs_set_eq:NN \CJKunderline \uline
        \cs_set_eq:NN \CJKunderdblline \uuline
        \cs_set_eq:NN \CJKunderwave \uwave
        \cs_set_eq:NN \CJKsout \sout
        \cs_set_eq:NN \CJKxout \xout
        \cs_set_eq:NN \CJKunderdot \dotuline
        \cs_set_eq:NN \CJKunderdashline \dashuline
      }
  }
  {
    \sys_if_engine_pdftex:TF
      {
        \cs_new_protected:Npn \hit@engine@postload:
          {
            \cs_set_eq:NN \CJKunderline \uline
            \cs_set_eq:NN \CJKunderdblline \uuline
            \cs_set_eq:NN \CJKunderwave \uwave
            \cs_set_eq:NN \CJKsout \sout
            \cs_set_eq:NN \CJKxout \xout
            \cs_set_eq:NN \CJKunderdot \dotuline
            \cs_set_eq:NN \CJKunderdashline \dashuline
          }
      }
      { \cs_new_protected:Npn \hit@engine@postload: { } }
  }
%    \end{macrocode}
%
% \textbf{三、下划线的位置与粗细按 Word。}Word 是一个 run 一条线，线心在
% 基线下 0.135\,em、粗 0.05\,em，纯汉字、纯西文、数字、混排四种 run 一个数。
% 自己造 docx 在 macOS Word 上量的（中易宋体 + Times，12/24/48\,pt 三档；
% 大字号躲开 Word 输出坐标 1/300 英寸的取整，48\,pt 那档四种 run 全落在
% 0.135/0.050 上，12\,pt 那档在 0.123～0.173 之间晃，那是它按设备像素取整）。
% 与 iota-hit 的实测同数。粗体西文 run 在 Word 里线粗一倍（0.10\,em），
% 这一档不另设。
%
% 不设的话 \pkg{ulem} 按字体的 (j 深度自动取，\pkg{xeCJKfntef} 默认 0.20\,em，
% 两边都在基线下 0.235\,em 上下、粗 0.4\,pt（12\,pt 上 0.033\,em），比 Word
% 低了近一倍、细了三成。
%
% \pkg{ulem} 的线由 \cs{ULdepth}（下缘）与 \cs{ULthickness} 定：下缘取
% 0.16\,em、粗 0.05\,em，线心正好 0.135。\cs{ULdepth} 是长度寄存器，不能在
% 加载时按 em 定死，改成在它自动取深度那一步（\cs{UL@setULdepth}）按当时
% 字体算。\pkg{xeCJKfntef} 的 \texttt{depth} 是线的\emph{上缘}（实测 0.16 给
% 出线心 0.185），取 0.11\,em；那对键收的是 token 列表，用时才求值，直接写
% em。两边都量过：12\,pt 上线心 0.1350、粗 0.0500。挂在两个包的加载钩子上，
% 哪个引擎装哪个都吃得到。
%    \begin{macrocode}
\AddToHook { package / ulem / after }
  {
    \cs_set:Npn \UL@setULdepth { \dim_set:Nn \ULdepth { 0.16em } }
    \cs_set:Npn \ULthickness { 0.05em }
  }
\AddToHook { package / xeCJKfntef / after }
  { \xeCJKsetup { underline = { depth = 0.11em , thickness = 0.05em } } }
%    \end{macrocode}
%
% 伪粗体不分引擎：系数是按 Word 实测定的数字（见 \file{src/engine/hit-fonts.dtx}），
% 数字两个引擎都吃（\pkg{luaotfload} 的 \texttt{embolden} 只是不吃布尔值）。
%
%    \begin{macrocode}
%</shared-engine>
%    \end{macrocode}
%
% \Finale
